\usepackage{mathrsfs}
\usepackage[all,cmtip]{xy}
\let\labelstyle\textstyle
\let\epsilon\varepsilon
\newcommand\mto{\mathchoice{\longmapsto}{\mapsto}{\mapsto}{\mapsto}}
\newcommand\hto{\mathrel{\lhook\joinrel\to}}
\newcommand{\tto}{\mathchoice{\to\hspace*{-5mm}\to}{\to\hspace*{-2.85mm}\to}{\to\hspace*{-3mm}\to}{\to\hspace*{-3mm}\to}}
\newcommand\iso{\stackrel{\sim}{\longrightarrow}}
\def\To#1{\mathchoice{\xrightarrow{\textstyle\kern4pt#1\kern3pt}}{\stackrel{#1}{\longrightarrow}}{}{}}
\def\Mto#1{\mathrel{\mapstochar\To{#1}}}
\def\Hto#1{\mathrel{\lhook\joinrel\To{#1}}}
\newcommand{\Changel}{\mathcode`l="7160}
\newcommand{\Changelback}{\mathcode`l="716C}
\newcommand{\NoChangel}[1]{%
\expandafter\let\csname old\string#1\endcsname=#1
\let#1=\relax
\newcommand{#1}{\mathcode`l="716C\csname old\string#1\endcsname\mathcode`l="7160 }%
}

\NoChangel{\log}
\NoChangel{\ln}
\NoChangel{\lim}
\NoChangel{\limsup}
\NoChangel{\liminf}
\NoChangel{\varlimsup}
\NoChangel{\varliminf}
\NoChangel{\varinjlim}
\NoChangel{\varprojlim}

\newcommand{\divv}{\mathrm{div}}
\newcommand{\fib}{\mathrm{fib}}
\let\ts\textstyle

\usepackage{tikz}
\usetikzlibrary{cd,arrows,positioning}
\tikzset{>=latex}
\tikzcdset{arrow style=Latin Modern}
\tikzset{link/.style={column sep=1.8cm,row sep=0.25cm}}
\tikzset{map/.style={row sep=0em, column sep=0em}}

\newcommand{\dradpl}{{\hspace*{-2mm}\begin{tikzcd}[ampersand replacement=\&,column sep = scriptsize]\mbox{}\ar[r,dashed]\&\mbox{}\end{tikzcd}\hspace*{-2mm}}}
\newcommand{\dratst}{{\hspace*{-2mm}\begin{tikzcd}[ampersand replacement=\&,column sep = small]\mbox{}\ar[r,dashed]\&\mbox{}\end{tikzcd}\hspace*{-2mm}}}
\let\drascst\dratst
\newcommand{\dra}{\mathchoice{\dradpl}{\dratst}{\drascst}{\drascst}}
\newcommand{\ps}{\!\!\xymatrix@C=3.5mm{\ar@{.>}[r]&}\!\!}

\DeclareMathOperator{\tr}{tr}
\DeclareMathOperator{\Bir}{Bir}
\DeclareMathOperator{\Spec}{Spec}
\DeclareMathOperator{\rk}{rk}
\DeclareMathOperator{\Pic}{Pic}
\DeclareMathOperator{\Lef}{Lef}
\DeclareMathOperator{\Fix}{Fix}
\DeclareMathOperator{\Aut}{Aut}
\DeclareMathOperator{\Bl}{Bl}
\DeclareMathOperator{\pr}{pr}
\DeclareMathOperator{\Ker}{Ker}
\DeclareMathOperator{\PGL}{PGL}
\DeclareMathOperator{\GL}{GL}
\DeclareMathOperator{\SL}{SL}
\DeclareMathOperator{\PSL}{PSL}
\DeclareMathOperator{\SO}{SO}
\DeclareMathOperator{\inv}{\boldsymbol{g}}
\DeclareMathOperator{\Nef}{Nef}
\DeclareMathOperator{\Mov}{Mov}
\DeclareMathOperator{\dP}{dP}
\DeclareMathOperator{\stab}{stab}
\DeclareMathOperator{\Supp}{Supp}
\DeclareMathOperator{\Exc}{Exc}
\DeclareMathOperator{\Bs}{Bs}
\DeclareMathOperator{\BirMori}{BirMori}

\theoremstyle{plain}
\newtheorem{thm}{Theorem}[section]
\newtheorem{lem}[thm]{Lemma}
\newtheorem{cor}[thm]{Corollary}
\newtheorem{prop}[thm]{Proposition}

\newtheorem{maintheorem}{Theorem}
\renewcommand\themaintheorem{\Alph{maintheorem}}
\newtheorem{maincorollary}[maintheorem]{Corollary}

\theoremstyle{definition}
\newtheorem{mydef}[thm]{Definition}
\newtheorem{rem}[thm]{Remark}
\newtheorem*{question}{Question}
\newtheorem{ex}[thm]{Example}

\newcommand{\I}{\ensuremath{\mathrm{I}}}
\newcommand{\II}{\ensuremath{\mathrm{II}}}
\newcommand{\III}{\ensuremath{\mathrm{III}}}
\newcommand{\IV}{\ensuremath{\mathrm{IV}}}
\newcommand{\TypeI}{\II_3/\pt}
\newcommand{\TypeII}{\II_3/\PP^1}

\newcommand{\CC}{\mathbb C}
\newcommand{\QQ}{\mathbb Q}
\newcommand{\FF}{\mathbb F}
\newcommand{\id}{\mathrm{id}}
\newcommand{\NN}{\mathbb N}
\newcommand{\kk}{\Bbbk}
\newcommand{\RR}{\mathbb R}
\newcommand{\PP}{\mathbb P}
\newcommand{\ZZ}{\mathbb Z}
\newcommand{\pt}{\mathrm{pt}}
\newcommand{\DD}{\mathcal D}
\newcommand{\XC}{X_{\mathbb C}}
\newcommand{\XR}{X_{\mathbb R}}
\newcommand{\MM}{\mathcal M}
\newcommand{\Bertini}{\mathscr B}
\newcommand*\conj[1]{\overline{#1}}

